% Dr. Rishi Gurnani — Resume v5
% Target: NVIDIA — AI for Chemistry & Materials, Developer Relations / Research Partnerships
% Compile with XeLaTeX or LuaLaTeX on Overleaf

\documentclass[10pt,letterpaper]{article}

\usepackage[top=0.65in, bottom=0.65in, left=0.75in, right=0.75in]{geometry}
\usepackage{fontspec}
\setmainfont[
  BoldFont   = TeX Gyre Heros Bold,
  ItalicFont = TeX Gyre Heros Italic
]{TeX Gyre Heros}
\usepackage{microtype}
\usepackage{parskip}
\setlength{\parskip}{0pt}
\usepackage{enumitem}
\setlist[itemize]{leftmargin=1.2em, itemsep=3pt, parsep=0pt, topsep=3pt}
\usepackage{tabularx}
\usepackage{array}
\usepackage{hyperref}
\hypersetup{colorlinks=false, hidelinks}
\usepackage{fontawesome5}
\usepackage{titlesec}

\titleformat{\section}{\bfseries\small\scshape\uppercase}{}{0pt}{}
  [\vspace{2pt}\hrule\vspace{4pt}]
\titlespacing*{\section}{0pt}{10pt}{4pt}

\newcommand{\role}[3]{%
  \noindent\textbf{#1}\hfill{\small #2}\\
  {\small\textit{#3}}\par\vspace{3pt}%
}

% Boxed hyperlink — thin rule, tight padding
\newcommand{\lbox}[2]{%
  {\setlength{\fboxsep}{1.5pt}\fbox{\href{#1}{#2}}}%
}

\begin{document}
\pagestyle{empty}

% ── Header ───────────────────────────────────────────────────────────────────
\begin{center}
  {\fontsize{20pt}{24pt}\selectfont\bfseries DR.\ RISHI GURNANI}\\[5pt]
  {\small
    U.S.\ Citizen\quad|\quad
    \lbox{mailto:rishipgurnani@gmail.com}{rishipgurnani@gmail.com}\quad|\quad
    \lbox{https://rishigurnani.github.io}{rishigurnani.github.io}\quad|\quad
    \lbox{https://github.com/rishigurnani}{\faGithub\ rishigurnani}\quad|\quad
    $>$1,900 citations\quad
  }
\end{center}
{\rule{\linewidth}{1.2pt}}
\vspace{4pt}

% ── Summary ──────────────────────────────────────────────────────────────────
\section{Summary}
\textbf{Brief:} 12+ years experience in materials chemistry; 10+ years experience in computer programming; 8+ years experience in deep learning; 8+ experience in research; 2+ years experience as a data science people leader \\\\
\textbf{Long:} AI researcher and engineer specializing in generative modeling and deep learning for chemistry
and materials science. Ph.D.\ from Georgia Tech; 16 publications ($>$1,900 citations, h-index 15),
including first-author papers in \textit{Nature Communications} and \textit{Chemistry of Materials}.
Proven ability to communicate complex AI methods to both technical and non-technical audiences —
as a people leader, keynote speaker, podcast guest, and invited lecturer at universities and national labs.
Deep familiarity with the academic and research ecosystem in AI for materials discovery.

% ── Core Competencies ────────────────────────────────────────────────────────
\section{Core Competencies}

\renewcommand{\arraystretch}{1.3}
\begin{tabularx}{\linewidth}{@{} >{\bfseries\small}p{0.21\linewidth} X @{}}
  Generative AI       & VAE, GAN, graph-based generative models; applied to molecular and polymer design \\
  Deep Learning       & GNNs, multitask learning, physics-informed ML, model ensembling, inference optimization \\
  Statistical         & Regression, classification, uncertainty quantification, feature engineering \\
  Research Ecosystem  & Active knowledge of leading labs, researchers, and open problems in AI for chemistry \& materials \\
  Programming         & Python, PyTorch, PyTorch Geometric, NumPy, Scikit-learn, Pandas, FastAPI, HuggingFace, PyTest \\
  Open Source         & \lbox{https://github.com/rishigurnani/polygnn_trainer}{\texttt{polygnn\_trainer}} (PyTorch ML training library, adopted by materials informatics community) \\
\end{tabularx}

% ── Work Experience ───────────────────────────────────────────────────────────
\section{Work Experience}

\role{Senior Manager, Data Science for Global Ingredient Research — Coca-Cola}{Feb 2025 – present}{Atlanta, GA}
\begin{itemize}
  \item Developing AI strategy and writing core algorithms for global R\&D in the following areas: sweetener application, functional ingredients for wellness, flavors, emulsifiers
  \item Algorithm stack: Sentence Transformers, LLMs, PLS Regression, Rule-based recommendation systems
  \item Tech Stack: OpenAI, PyTorch, Scikit-Learn, FastAPI, HuggingFace
\end{itemize}

\vspace{5pt}

\role{Director of Software Engineering \& Algorithms — Matmerize Inc.}{Mar 2021 – Feb 2025}{Atlanta, GA}
\vspace{3pt}
\begin{itemize}
  \item Developed the deep learning core of PolymRize\texttrademark, an enterprise ML platform
    ($\sim$100K lines of code, 5 languages) adopted by $>$20 research organizations and companies.

  \item Served as primary technical point of contact for academic and industry partners,
    translating research needs into software capabilities and communicating algorithmic methods
    to audiences ranging from wet-lab chemists to ML engineers.

  \item Mentored junior developers and interns; alumni now at \textbf{Meta}, 
    \textbf{UC Berkeley (PhD)}, and \textbf{Georgia Tech (PhD)}.
\end{itemize}

% ── Education ─────────────────────────────────────────────────────────────────
\newpage
\section{Education}

\begin{tabularx}{\linewidth}{@{} X r @{}}
  \textbf{Ph.D., Materials Science \& Engineering} — Georgia Tech (GPA 4.00) & 2019–2023 \\
  \multicolumn{2}{@{}l@{}}{\small\textit{Dissertation: Advances in AI for Accelerated Discovery of Energy Storage Polymers (GNNs \& Generative ML)}}\\[3pt]
  \textbf{M.S., Computer Science} — Georgia Tech (Specialization: Machine Learning) & \\[2pt]
  \textbf{B.S., Materials Science \& Engineering} — UC Berkeley & 2014–2018 \\
\end{tabularx}

% ── Awards & Talks ────────────────────────────────────────────────────────────
\section{Awards, Talks \& Open Source}

\begin{tabularx}{\linewidth}{@{} X r @{}}
  \textbf{Best Ph.D.\ Thesis}, Georgia Tech Sigma Xi        & 2024 \\
  \textbf{Graduate Student Award}, Materials Research Society & 2023 \\
  \textbf{Gold Medal, Best Presentation}, 2nd Energy \& Informatics International Forum & 2022 \\
\end{tabularx}

\vspace{5pt}

\noindent\textbf{Invited Talks:} Keynote, \lbox{https://www.riseconf.net/}{2024 RISE Conference};
Tokyo Institute of Technology (2022); Sandia National Laboratory (2021); nanoHUB ML Series (2021);
STLE 76th Annual Meeting (2022). Podcast: \lbox{https://www.youtube.com/watch?v=srLJlK67tXc}{\textit{It's a Material World}} (2023).

\vspace{5pt}

\noindent\textbf{Open Source:}
\lbox{https://github.com/rishigurnani/polygnn_trainer}{\faGithub\ \texttt{polygnn\_trainer}} —
PyTorch library for automated ML training, adopted by the materials informatics community.\quad
\lbox{https://github.com/rishigurnani/fintracker}{\faGithub\ \texttt{fintracker}} —
open-source personal finance analytics library.

% ── Publications ─────────────────────────────────────────────────────────────
\section{Publications \ {\normalfont\small(16 papers, $>$1,900 citations, h-index 15)}}

\begin{enumerate}[leftmargin=1.4em, itemsep=4pt, parsep=0pt, topsep=3pt]

  \item \lbox{https://www.nature.com/articles/s41467-024-50413-x}{AI-assisted discovery of high-temperature dielectrics for energy storage.}
    R Gurnani, S Shukla et al. \textit{Nature Communications} 15 (1), 2024. \textit{(93 citations)}

  \item \lbox{https://pubs.acs.org/doi/full/10.1021/acs.chemmater.2c02991}{Polymer informatics at scale with multitask graph neural networks.}
    R Gurnani, C Kuenneth, A Toland, R Ramprasad. \textit{Chemistry of Materials} 35 (4), 2023. \textit{(106 citations)}

  \item \lbox{https://rishigurnani.files.wordpress.com/2021/08/polyg2g-a-novel-machine-learning-algorithm-applied-to-the-generative-design-of-polymer-dielectrics.pdf}{PolyG2G: A novel machine learning algorithm applied to the generative design of polymer dielectrics.}
    R Gurnani, D Kamal, H Tran et al. \textit{Chemistry of Materials} 33 (17), 2021. \textit{(82 citations)}

  \item \lbox{https://www.nature.com/articles/s41578-024-00708-8}{Design of functional and sustainable polymers assisted by artificial intelligence.}
    H Tran, R Gurnani, C Kim et al. \textit{Nature Reviews Materials} 9 (12), 2024. \textit{(198 citations)}

  \item Stereoisomerism of Vicinal Polydichloronorbornene for Ultra-High-Temperature Capacitive Energy Storage.
    J Hao, S Shukla, R Gurnani et al. \textit{Advanced Materials} 37 (15), 2025.

  \item A physics-enforced neural network to predict polymer melt viscosity.
    A Jain, R Gurnani, A Rajan, HJ Qi, R Ramprasad. \textit{npj Computational Materials} 11 (1), 2025.

  \item Rationally designed high-temperature polymer dielectrics for capacitive energy storage: An experimental and computational alliance.
    PS Aklujkar, R Gurnani et al. \textit{Progress in Polymer Science} 161, 2025.

  \item Polymer composites informatics for flammability, thermal, mechanical and electrical property predictions.
    H Tran, C Kim, R Gurnani et al. \textit{Polymer Chemistry} 16 (30), 2025.

  \item Gas permeability, diffusivity, and solubility in polymers: Simulation-experiment data fusion and multi-task machine learning.
    BK Phan, KH Shen, R Gurnani et al. \textit{npj Computational Materials} 10 (1), 2024.

  \item Pendant group functionalization of cyclic olefin for high temperature and high-density energy storage.
    S Shukla, C Wu, A Mishra, R Gurnani et al. \textit{Advanced Materials} 36 (52), 2024.

  \item Effect of fluorine in redesigning energy-storage properties of high-temperature dielectric polymers.
    AA Deshmukh, C Wu, R Gurnani et al. \textit{ACS Applied Materials \& Interfaces} 15 (40), 2023.

  \item Identifying High-Performance Metal-Organic Frameworks for Low-Temperature Oxygen Recovery from Helium by Computational Screening.
    S Jamdade, R Gurnani, H Fang, SE Boulfelfel, R Ramprasad, DS Sholl. \textit{Industrial \& Engineering Chemistry Research} 62 (4), 2023.

  \item Interpretable machine learning-based predictions of methane uptake isotherms in metal-organic frameworks.
    R Gurnani, Z Yu, C Kim, DS Sholl, R Ramprasad. \textit{Chemistry of Materials} 33 (10), 2021.

  \item \lbox{https://rishigurnani.wordpress.com/wp-content/uploads/2020/11/5.0023759.pdf}{Machine-learning predictions of polymer properties with Polymer Genome.}
    H Doan Tran, C Kim, L Chen, R Gurnani et al. \textit{Journal of Applied Physics} 128 (17), 2020. \textit{(320 citations)}

  \item Experimental characterization of the shear properties of 3D-printed ABS and polycarbonate parts.
    S Rohde, J Cantrell, R Gurnani et al. \textit{Experimental Mechanics} 58 (6), 2018.

  \item Experimental characterization of the mechanical properties of 3D-printed ABS and polycarbonate parts.
    JT Cantrell, S Rohde, D Damiani, R Gurnani et al. \textit{Rapid Prototyping Journal} 23 (4), 2017. \textit{(777 citations)}

\end{enumerate}

\end{document}